\section{Startanleitung}
\label{sec:startanleitung}

Dieser Abschnitt beschreibt, wie das TaskGrid+{}-System gestartet, betrieben
und wieder gestoppt wird. Alle Schritte setzen voraus, dass Docker und
Docker Compose auf dem System installiert sind.

\subsection{Voraussetzungen}

Folgende Software muss auf dem Rechner vorhanden sein:

\begin{itemize}
  \item \textbf{Docker} (Version 24 oder höher empfohlen)
  \item \textbf{Docker Compose} (als Plugin: \texttt{docker compose}, oder als
        eigenständiges Binary \texttt{docker-compose})
  \item Das Repository muss lokal ausgecheckt sein
\end{itemize}

Ob Docker korrekt installiert ist, lässt sich mit folgendem Befehl prüfen:

\begin{lstlisting}[language=bash,caption={Docker-Versionsprüfung}]
docker --version
docker compose version
\end{lstlisting}

\subsection{System starten}

Aus dem Wurzelverzeichnis des Repositories genügt ein einziger Befehl,
um alle Komponenten zu bauen und zu starten:

\begin{lstlisting}[language=bash,caption={System starten}]
docker compose up --build
\end{lstlisting}

Die Option \texttt{--build} stellt sicher, dass alle Docker-Images neu gebaut werden.
Beim ersten Start dauert das je nach Internetverbindung einige Minuten, da die
Python-Basisimages geladen werden. Folgeaufrufe ohne Codeänderungen sind deutlich schneller.

Um das System im Hintergrund zu starten (detached mode), wird \texttt{-d} ergänzt:

\begin{lstlisting}[language=bash,caption={System im Hintergrund starten}]
docker compose up --build -d
\end{lstlisting}

\subsubsection*{Startreihenfolge}

Docker Compose startet die Komponenten in der richtigen Reihenfolge,
da jeder Service Health-Checks und \texttt{depends\_on}-Bedingungen definiert:

\begin{enumerate}
  \item \textbf{namensdienst} startet zuerst und öffnet seinen gRPC-Port (50052).
  \item \textbf{dispatcher} startet, sobald der Namensdienst als \textit{healthy}
        gemeldet wird, und öffnet Port 50051 (gRPC) sowie 8080 (HTTP-Monitoring).
  \item \textbf{worker-*} starten, sobald Namensdienst und Dispatcher bereit sind.
        Jeder Worker registriert sich beim Namensdienst und beginnt mit dem Senden
        von Heartbeats.
  \item \textbf{client} startet als letzter, führt eine Demo aus und beendet sich.
\end{enumerate}

\subsection{Worker skalieren}

Einer der Hauptvorteile der Architektur ist die einfache horizontale Skalierbarkeit.
Einzelne Worker-Typen lassen sich mit dem \texttt{--scale}-Flag hochskalieren,
ohne dass am Dispatcher oder Namensdienst etwas geändert werden muss:

\begin{lstlisting}[language=bash,caption={Mehrere Worker desselben Typs starten}]
docker compose up --build -d --scale worker-sum=3
\end{lstlisting}

Auf diese Weise starten drei Instanzen des \texttt{worker-sum}-Containers gleichzeitig.
Jede Instanz registriert sich mit einer eigenen eindeutigen ID beim Namensdienst.
Der Dispatcher verteilt eingehende Tasks dann automatisch per Round-Robin auf
alle drei Instanzen. Mehrere Typen lassen sich kombinieren:

\begin{lstlisting}[language=bash,caption={Verschiedene Worker-Typen gleichzeitig skalieren}]
docker compose up --build -d --scale worker-sum=3 --scale worker-reverse=2
\end{lstlisting}

\subsection{System überprüfen}

Sobald das System läuft, können die einzelnen Komponenten auf verschiedene Arten
überprüft werden.

\subsubsection*{Container-Status}

Den aktuellen Status aller laufenden Container zeigt:

\begin{lstlisting}[language=bash,caption={Laufende Container anzeigen}]
docker compose ps
\end{lstlisting}

\subsubsection*{Logs einsehen}

Die Log-Ausgaben einzelner Komponenten lassen sich gezielt abrufen:

\begin{lstlisting}[language=bash,caption={Logs einer Komponente anzeigen}]
docker logs dispatcher
docker logs namensdienst
docker compose logs worker-sum
\end{lstlisting}

\subsubsection*{Monitoring-Endpoint}

Der Dispatcher stellt unter Port 8080 einen HTTP-Monitoring-Endpoint bereit,
der aktuelle Systemkennzahlen als JSON ausgibt:

\begin{lstlisting}[language=bash,caption={Systemstatus abfragen}]
curl http://localhost:8080/status
\end{lstlisting}

Die Ausgabe enthält unter anderem die Anzahl aktiver Worker, offene und laufende
Tasks sowie Statistiken zu Timeouts und Retries. Ein Beispiel:

\begin{lstlisting}[language={},caption={Beispiel-Ausgabe des Monitoring-Endpoints}]
{
  "aktive_worker": 3,
  "unterstützte_tasktypen": ["sum", "reverse", "hash", "upper", "wait"],
  "offene_tasks": 0,
  "laufende_tasks": 1,
  "abgeschlossene_tasks": 12,
  "fehlgeschlagene_tasks": 0,
  "durchschnittliche_bearbeitungszeit_ms": 42.5,
  "anzahl_timeouts": 0,
  "anzahl_retries": 0
}
\end{lstlisting}

\subsubsection*{Einen Task manuell einreichen}

Tasks können über den mitgelieferten Client-Container oder direkt per gRPC
an den Dispatcher gesendet werden. Der Client bietet zwei Modi:

\begin{lstlisting}[language=bash,caption={Task senden und Ergebnis abfragen}]
# Task einreichen (gibt eine task_id zurueck)
docker compose run --rm client python src/client/client.py send sum "1,2,3,4"

# Ergebnis einer Task-ID abfragen
docker compose run --rm client python src/client/client.py result 1
\end{lstlisting}

\subsection{System stoppen}

Zum sauberen Beenden aller Container genügt:

\begin{lstlisting}[language=bash,caption={System stoppen}]
docker compose down
\end{lstlisting}

Sollen zusätzlich die persistierten Log-Volumes entfernt werden, wird
\texttt{-v} ergänzt:

\begin{lstlisting}[language=bash,caption={System stoppen und Volumes löschen}]
docker compose down -v
\end{lstlisting}